% =======================================================================================
% Chapter 3 -- Methodology
% SECTION-H(12): all materials and data explained, supported with figures and tables;
% analytical techniques and experiments described in detail.
% Chapters 3-5 together must not exceed forty pages.
% =======================================================================================
\chapter{Methodology}
\label{ch:methodology}

\section{Research Design}
\label{sec:design}

The study follows a quantitative, experimental design within a positivist frame: a working system
is constructed, and its behaviour is characterised by instrumental measurement on data it has never
seen. Four experiments are reported.

\begin{enumerate}[label=\textbf{E\arabic*.}]
  \item \textbf{Synthesis experiment.} The complete cascade is run over the held-out test split in
        both directions, and recognition accuracy, translation quality, speaker identity,
        intelligibility and prosody are measured for every utterance
        (Sections~\ref{sec:architecture}--\ref{sec:instruments}).
  \item \textbf{Corpus study.} The cross-lingual prosodic relationship between English and Spanish
        is characterised for the same human speaker over the whole corpus
        (Section~\ref{sec:corpus_study}).
  \item \textbf{Prosody prediction experiment.} A supervised mapping from source-language to
        target-language prosody is fitted on the training split, selected on development data and
        tested once, against three baselines (Section~\ref{sec:ridge}).
  \item \textbf{Fine-tuning experiment.} The translation component is adapted to the in-domain
        conversational data and compared against its pre-trained baseline under identical decoding
        (Section~\ref{sec:finetune_method}).
\end{enumerate}

A fifth deliverable, the dubbing application (Section~\ref{sec:webapp}), is engineering rather than
experiment; it realises Objective~3 and is described because it is a research output and because
the retiming stage it contains follows directly from a measured result.

The design has one defining property. Because every test utterance has a matching human recording
of \emph{the same speaker performing the same content in the target language}, each measurement is
made against a human reference rather than in the abstract. That is what converts an uninterpretable
similarity score into a percentage of an achievable ceiling.

\section{Materials: the DRAL Corpus}
\label{sec:corpus}

All experimental data come from the 16\,kHz release of the Dialogs Re-enacted Across Languages
corpus \cite{ward2023dral}, obtained from the University of Texas at El Paso and placed in the
public domain by its authors. In the collection protocol, bilingual participants hold a
spontaneous conversation in one language; the recording is then segmented, and the same
participants re-enact each fragment in the other language, aiming to reproduce not only the content
but the manner of the original. Both sides of a pair therefore come from one person performing one
communicative act twice.

Only the \texttt{fragments-short} portion of the release was used. The \texttt{fragments-long} and
\texttt{recordings} portions contain longer or unsegmented versions of the same material and carry
no additional pair structure, and the Japanese material in the release is out of scope. Pairs were
resolved through the official \texttt{trans\_id} links in
\pth{metadata/fragments-short.csv} rather than by substituting the language prefix in a filename,
because a number of valid pairs use different conversation or channel identifiers. Human
transcripts were merged from the official \pth{EN_transcripts.txt} and
\pth{ES_transcripts.txt} files; where a transcript is absent the field is left empty and the
affected metric is reported as unavailable, never imputed from recogniser output.
Table~\ref{tab:corpus} records the resulting inventory, which was validated locally rather than
taken from the documentation.

\begin{table}[htbp]
\caption{Inventory of the DRAL 16\,kHz \texttt{fragments-short} release as validated locally. All
figures were produced by \texttt{bvt inspect-dral} and \texttt{bvt build-manifest}.}
\label{tab:corpus}
\centering
{\singlespacing\small
\begin{tabular}{@{}lr@{}}
\toprule
\textbf{Property} & \textbf{Value} \\
\midrule
Linked English--Spanish fragment pairs & 2{,}893 \\
Audio files opened successfully        & 5{,}786 \\
Total speech duration                  & 4.27 hours \\
Files at 16\,kHz                       & 5{,}784 \\
Unique speaker identifiers             & 96 \\
Pairs performed by one speaker         & 2{,}834 \\
Pairs with differing linked speaker identifiers & 59 \\
Non-empty English human transcripts    & 2{,}871 \\
Non-empty Spanish human transcripts    & 2{,}881 \\
\bottomrule
\end{tabular}\par}
\end{table}

\subsection{Why This Corpus and Not the Ones Named in the Proposal}
\label{sec:corpus_choice}

The approved proposal named Mozilla Common Voice \cite{ardila2020commonvoice} and the CVSS
speech-to-speech corpus \cite{jia2022cvss}. Both were examined and set aside, for one reason. In
Common Voice the English and Spanish recordings come from different people, so there is no
same-speaker cross-language reference. In CVSS the target-language speech is itself synthesised by
a text-to-speech system, so comparing generated speech against it would compare one synthesiser
with another rather than with a human. Neither corpus can supply the human ceiling on which the
central measurement of this study depends. DRAL can, because one bilingual person performs both
sides. The cost of the choice is scale --- 4.27 hours against thousands --- and it is accepted
deliberately: this study measures rather than trains, and measurement needs a valid reference more
than it needs volume.

Two files, \texttt{EN\_072\_7.wav} and \texttt{ES\_072\_7.wav}, are stored at 44.1\,kHz despite the
name of the release; they are resampled at read time for feature extraction and speaker embedding,
and the raw data are not modified. One line of \pth{ES_transcripts.txt} lacks a field delimiter;
the parser recovers the annotator code and emits a warning rather than editing the source file.

\subsection{Speaker-Disjoint Partitioning}
\label{sec:splits}

The corpus was partitioned into training, development and test sets in the proportions
70/15/15 with random seed 498. A conventional random split would be invalid here, because a speaker
appearing in both training and test data would let the evaluation reward memorised voice
characteristics rather than generalisation. Assignment therefore operates on \emph{speakers}, not
on pairs.

Because 59 pairs link fragments whose English and Spanish sides carry different speaker
identifiers, the two identifiers of such a pair must never be separated. The splitter first builds
a graph in which an edge joins the two speaker identifiers of every pair, computes the connected
components of that graph, and assigns whole components to splits. This guarantees that no speaker,
and no linked group of speakers, appears on both sides of the train/test boundary; the property is
asserted programmatically after assignment. The resulting split contains 2{,}022 training pairs,
436 development pairs and 435 test pairs. The test split, on which all reported system results are
computed, contains 14 speakers.

\section{System Architecture}
\label{sec:architecture}

The system is a cascade of four pre-trained neural components, illustrated in
Figure~\ref{fig:architecture}. Three participate in generation; the fourth is used only for
measurement.

\begin{figure}[htbp]
\centering
\begin{tikzpicture}[
  box/.style={draw, rounded corners=2pt, minimum height=1.05cm, minimum width=2.1cm,
              align=center, font=\small, fill=blue!5, inner sep=2pt},
  measbox/.style={draw, rounded corners=2pt, minimum height=1.0cm, minimum width=2.6cm,
              align=center, font=\small, fill=orange!10, dashed, inner sep=2pt},
  data/.style={align=center, font=\scriptsize\itshape, inner sep=2pt},
  arr/.style={-{Latex[length=2mm]}, thick},
  measline/.style={-{Latex[length=1.8mm]}, semithick, dashed, gray!70},
  lbl/.style={font=\scriptsize, inner sep=1pt, fill=white}
]
% ---- generation path -----------------------------------------------------------------
\node[data]                        (src)  at (0,0)    {source\\audio};
\node[box,  right=0.7cm of src]    (asr)               {Whisper\\[-1pt]\scriptsize small};
\node[box,  right=0.75cm of asr]   (mt)                {MarianMT\\[-1pt]\scriptsize opus-mt};
\node[box,  right=1.25cm of mt]    (tts)               {XTTS-v2};
\node[data, right=0.7cm of tts]    (out)               {dubbed\\audio};

\draw[arr] (src) -- (asr);
\draw[arr] (asr) -- node[lbl, above] {text}        (mt);
\draw[arr] (mt)  -- node[lbl, above] {translation} (tts);
\draw[arr] (tts) -- (out);

% voice conditioning: the source waveform reaches the synthesiser, but only as a voice
% reference -- no pitch contour, no timing.
\draw[arr] (src.south) -- ++(0,-0.85)
           -- node[lbl, below, pos=0.55] {reference audio (voice only)} ($(tts.south)+(0,-0.85)$)
           -- (tts.south);

% ---- measurement layer ---------------------------------------------------------------
\node[measbox] (ecapa)   at ($(mt.south)+(-0.9,-3.3)$) {ECAPA-TDNN\\[-1pt]\scriptsize speaker embedding};
\node[measbox, right=0.9cm of ecapa] (pros) {pyin / RMS\\[-1pt]\scriptsize prosody features};
\node[data, left=0.9cm of ecapa] (human) {human target-language\\recording (DRAL)};

\draw[measline] (out.south)   -- ++(0,-1.55) -| (pros.north);
\draw[measline] ($(out.south)+(0,-1.55)$) -| (ecapa.north);
\draw[measline] (human.east)  -- (ecapa.west);
\draw[measline] (human.south) |- ($(pros.south)+(0,-0.55)$) -- (pros.south);
\draw[measline] (src.south)   -- ++(0,-0.4) -| ($(ecapa.north)+(-0.7,0)$);

\node[font=\scriptsize\itshape, gray!70, anchor=west] at ($(pros.east)+(0.15,0)$) {measurement only};
\end{tikzpicture}
\caption{Architecture of the cascaded English--Spanish speech-to-speech translation system. Solid
arrows are the generation path; the dashed lower layer is measurement only and never influences
generation. XTTS-v2 is conditioned on the source recording for voice, but is not given the source
pitch contour --- a design property that Section~\ref{sec:prosody_gap} shows to be consequential.}
\label{fig:architecture}
\end{figure}

\textbf{Recognition.} Whisper \texttt{small} \cite{radford2023whisper} transcribes the source
recording, with the source language supplied explicitly rather than auto-detected, so that a
recognition error cannot silently become a language-identification error.

\textbf{Translation.} MarianMT \cite{junczysdowmunt2018marian} in the two OPUS-MT bilingual
configurations \pth{Helsinki-NLP/opus-mt-en-es} and \pth{Helsinki-NLP/opus-mt-es-en}
\cite{tiedemann2020opusmt} translates the transcript, decoding with a maximum of 512 new tokens.

\textbf{Synthesis.} XTTS-v2 \cite{casanova2024xtts} synthesises the translation in the target
language, conditioned on the \emph{source} recording as the voice reference. The synthesiser
receives the target text and the reference waveform; it does not receive the source pitch contour,
its energy envelope or its timing.

\textbf{Measurement.} ECAPA-TDNN \cite{desplanques2020ecapa}, distributed as
\pth{speechbrain/spkrec-ecapa-voxceleb} \cite{ravanelli2021speechbrain}, produces speaker
embeddings. It generates nothing; it is an instrument.

Determinism is enforced per utterance. The synthesis seed is derived as the first four bytes of
\texttt{sha256(pair\_id\,|\,direction)} reduced modulo $2^{31}-1$, and the framework's random
number generator is seeded with it immediately before synthesis. Python's built-in \texttt{hash} is
salted per process and is deliberately not used. Every row of the results file records the seed
that produced it, so any single utterance can be regenerated exactly.

Failures are treated as data. If any stage raises, the row is written with
\texttt{status = failed}, the failing stage and the exception message, and it remains in the
results file and in the reported failure rate. No row is silently dropped.

\section{Evaluation Instruments}
\label{sec:instruments}

\subsection{Recognition Accuracy}
Word error rate is computed between the human transcript and the Whisper output using
\texttt{jiwer}, under a fixed normalisation applied identically to reference and hypothesis:
lower-casing, punctuation removal, whitespace collapsing and trimming. Both a corpus-level WER
(total edits divided by total reference words) and a mean over utterances are reported, because the
two answer different questions and diverge when short utterances are common. Utterances without a
human transcript are excluded from this metric only.

\subsection{Translation Quality}
Corpus BLEU \cite{papineni2002bleu} and chrF \cite{popovic2015chrf} are computed with
\texttt{sacrebleu} \cite{post2018sacrebleu} against the human transcript of the paired
target-language recording. That reference deserves a caution which is repeated wherever the
resulting scores appear: the DRAL target side is a \emph{re-enactment}, not a literal translation,
so the reference is a legitimate rendition of the same communicative content but not a
sentence-level gloss. The absolute scores are therefore lower than the same models achieve on
written parallel corpora, and only differences measured against the same reference are meaningful.

\subsection{Speaker Identity and the Human Ceiling}
\label{sec:identity_method}
Speaker similarity is the cosine between ECAPA embeddings of two recordings, computed after
down-mixing to mono and resampling to 16\,kHz. Two quantities are computed for every utterance:

\begin{align}
S_{\text{system}} &= \cos\!\big(e(\text{source recording}),\; e(\text{synthesised output})\big),
\label{eq:ssystem}\\[2pt]
S_{\text{human}}  &= \cos\!\big(e(\text{source recording}),\;
                                e(\text{human target-language recording})\big),
\label{eq:shuman}
\end{align}

where $e(\cdot)$ denotes the ECAPA embedding. $S_{\text{human}}$ is the \emph{ceiling}: it is how
similar the same person sounds to themselves when they switch language, measured on real human
recordings. Identity preservation is reported as the ratio
$100\,S_{\text{system}}/S_{\text{human}}$. This anchor is available only because DRAL is
same-speaker parallel; in a corpus without that property, $S_{\text{system}}$ would have no scale.

Both quantities are averaged first within each speaker and then across speakers. The test split is
markedly unbalanced --- one speaker contributes 118 pairs and another two --- so a plain mean over
rows would weight the result towards whoever happened to be recorded most. The speaker-level mean
gives every speaker equal weight and is the statistic reported throughout; row-level means are
given alongside it in Appendix~\ref{app:speakerlevel} for completeness.

\subsection{Intelligibility of Synthesised Speech}
The synthesised output is re-transcribed with Whisper and compared against the text that was sent
to the synthesiser. This isolates the synthesiser: a high value means the generated waveform is
hard to recognise, independently of whether the translation was any good. Rows are filtered on the
\emph{reference} being present, never on the hypothesis. An empty re-transcription means the audio
was unintelligible, which is a total failure scoring 1.0, not a missing observation; excluding such
rows would make the corpus figure look better than it is.

\subsection{Prosodic Features}
\label{sec:prosody_features}
Four acoustic descriptors are extracted from every recording with librosa
\cite{mcfee2015librosa}: fundamental frequency by the probabilistic pYIN estimator
\cite{mauch2014pyin}, frame-level root-mean-square energy, total duration, and speaking rate as
words per second from the corresponding text.

Two distinct frequency bands are used, and keeping them distinct is essential to the correctness of
the result.

\begin{itemize}
  \item The \textbf{search band}, 65--1000\,Hz, is the range given to pYIN. It must be wide.
        pYIN's voiced/unvoiced decision depends on the number of frequency candidates available,
        so narrowing the band lowers the aggregate voiced probability and the estimator stops
        detecting voicing at all on quiet recordings. This was measured, not assumed
        (Section~\ref{sec:f0correction}).
  \item The \textbf{plausibility band}, 60--400\,Hz, is applied downstream as a data-quality gate.
        A wide search band admits occasional octave doublings; a median F0 outside the plausible
        adult speaking range is a tracking failure rather than a voice, and such observations are
        excluded before any statistic is computed.
\end{itemize}

All F0 statistics are computed in \textbf{semitones} relative to a 100\,Hz reference,
\begin{equation}
s = 12\log_2\!\left(\frac{f}{100}\right),
\label{eq:semitones}
\end{equation}
where $f$ is the measured frequency in hertz, because pitch perception is logarithmic. A correlation computed in hertz is
dominated by whichever values lie furthest from the mean, and a single octave-doubled observation
can destroy it; the consequences of getting this wrong are reported in
Section~\ref{sec:f0correction}. Throughout, \emph{F0 level} denotes the mean semitone value of an
utterance and \emph{F0 range} its within-utterance standard deviation in semitones.

\subsection{Metrics Deliberately Not Used}
\label{sec:not_used}
The approved proposal named PESQ and STOI. Both are intrusive metrics: they compare a degraded
signal against a time-aligned clean recording of \emph{the same utterance}. In cross-lingual
synthesis no such reference can exist, because the output is in a different language and contains
different words of different durations. Reporting them would mean either fabricating an alignment
or comparing unrelated signals. They are therefore not reported, and the reason is stated rather
than left as an omission. A mean opinion score study was also planned; it was not conducted within
the available time and is recorded as a limitation in Section~\ref{sec:limitations} and as the
first item of future work.

\section{Corpus Study: Cross-Lingual Prosody in Human Speech}
\label{sec:corpus_study}

The corpus study characterises the relationship between the English and Spanish prosody of the
\emph{same} speaker, over the whole corpus rather than the test split, and supplies the human
reference against which the system is judged.

\subsection{Quality Gate}
Both sides of a pair must yield a usable measurement, so a pair is retained only if, on
\emph{both} sides: F0 was estimated at all; the voiced-frame ratio is at least 0.30; duration lies
between 0.4 and 20\,s; and mean F0 lies inside the plausibility band. Analysis is further
restricted to pairs performed by a single speaker, since the research question is about one
person's behaviour in two languages. The gate is a criterion on the \emph{measurement}, applied
before any outcome is computed, and every count it removes is reported in
Table~\ref{tab:gate}.

\subsection{Within- and Between-Speaker Decomposition}
\label{sec:decomposition}
A Pearson correlation computed over utterances pooled from many speakers conflates two effects. If
low-pitched people are low-pitched in both languages, a strong pooled correlation follows from
speaker identity alone, with no prosody transfer whatsoever. Every correlation in this dissertation
is therefore reported in three forms:

\begin{itemize}
  \item the \textbf{pooled} correlation over all utterances;
  \item the \textbf{between-speaker} correlation, computed over one point per speaker, namely
        their mean on each side;
  \item the \textbf{within-speaker} correlation, computed after centring both variables inside
        each speaker, which removes stable speaker differences and leaves only utterance-level
        variation.
\end{itemize}

The within-speaker figure is the one that answers the research question, and it is the figure
quoted in the discussion. The share of total variance attributable to speaker differences is
reported alongside, from a one-way decomposition
$\mathrm{SS}_{\text{between}}/(\mathrm{SS}_{\text{between}}+\mathrm{SS}_{\text{within}})$.

\section{Predicting Target-Language Prosody}
\label{sec:ridge}

The corpus study measures the correspondence; this experiment asks whether it can be
\emph{exploited}. For each direction and each of three targets --- target F0 level, target F0
range, and $\log(\text{target duration}/\text{source duration})$ --- a model predicts the target
value from six source-side features: F0 level, F0 range, log duration, voiced ratio, word count and
speaking rate.

\subsection{Baselines}
A model is only interesting if it beats what can be obtained without one. Three baselines are
fitted on the training split alone:

\begin{description}[leftmargin=2.6cm, style=nextline]
  \item[\textbf{B0}] \emph{copy source} --- predict the source value unchanged (and zero for the
        log duration ratio).
  \item[\textbf{B1}] \emph{train mean} --- predict the training-set mean of the target, ignoring
        the source entirely.
  \item[\textbf{B2}] \emph{copy plus offset} --- predict the source value plus a single global
        scalar, the mean residual on the training split. This is the baseline to beat: it captures
        any systematic language-level shift with one number.
\end{description}

\subsection{Ridge Regression on the Residual}
The model is ridge regression \cite{hoerl1970ridge} in closed form, with standardised inputs and
an unpenalised intercept, fitted to the \emph{residual from copy-source} rather than to the
absolute target:

\begin{equation}
\hat{\mathbf{w}} = \big(\mathbf{X}^{\!\top}\mathbf{X} + \lambda \mathbf{I}\big)^{-1}
                    \mathbf{X}^{\!\top}\big(\mathbf{y} - \mathbf{c}\big),
\qquad
\hat{y} = c + \mathbf{x}^{\!\top}\hat{\mathbf{w}} + b,
\label{eq:ridge}
\end{equation}

where $\mathbf{X}$ holds the standardised source features, $\mathbf{y}$ the targets, $\mathbf{c}$
the copy-source prediction, $\lambda$ the regularisation strength and $b$ the intercept. Fitting
the residual makes the model strictly \emph{nest} B2: as $\lambda \to \infty$ the weights vanish
and the prediction becomes copy-plus-mean-residual, which is exactly B2. The model can therefore
only improve on B2 up to estimation noise, which removes an entire class of misleading comparison.

$\lambda$ is selected on the development split from the grid
$\{10^{-3},10^{-2},\dots,10^{6}\}$ using the one-standard-error rule: the largest $\lambda$ whose
development error is within one standard error of the best. Plain minimisation is unstable on a
few-hundred-row development set, and because the model nests B2, an unlucky small $\lambda$ can
make it appear worse than the trivial baseline for reasons that are pure selection noise.
Preferring the most-shrunk admissible $\lambda$ biases towards the baseline, which is the
conservative choice. The test split is scored once.

\subsection{Statistical Inference with Speaker Clustering}
\label{sec:inference}
The test split holds 348 usable pairs (347 in \esto{}) but only 14 speakers, and errors are correlated
within a speaker. Treating each recording as an independent observation would overstate
significance. Two clustered procedures are reported instead of the naive per-recording test: a
Wilcoxon signed-rank test \cite{wilcoxon1945} on \emph{per-speaker mean absolute errors}
($n=14$), and a bootstrap confidence interval for the difference in mean absolute error that
resamples \emph{speakers} rather than recordings \cite{efron1993bootstrap}, with 2{,}000 draws at
seed 498. Because six models are fitted, the family-wise position is reported explicitly using
Bonferroni and Benjamini--Hochberg \cite{benjamini1995fdr} corrections rather than left implicit.

Two properties of the design produce results that look like defects and are not, so they are
stated in advance. For the log duration ratio, B1 and B2 are mathematically identical, since
copy-source is zero and copy-plus-offset reduces to the training mean; the comparison between them
is therefore exactly null. And the errors of B0, B1 and B2 are symmetric between directions by
construction, because exchanging source and target mirrors the error distribution; only the ridge
rows genuinely differ between the two directions.

\section{Fine-Tuning the Translation Component}
\label{sec:finetune_method}

The pre-trained translation models were trained on written text, whereas the corpus is spontaneous
bilingual conversation containing disfluencies, fillers and regional lexis. This experiment
measures what in-domain adaptation is worth. Each direction was fine-tuned independently from the
pre-trained checkpoint on the training split only, with the settings in
Table~\ref{tab:finetune_setup}.

\begin{table}[htbp]
\caption{Fine-tuning configuration. Decoding settings are identical for the pre-trained and
fine-tuned systems, so the comparison isolates the effect of adaptation.}
\label{tab:finetune_setup}
\centering
{\singlespacing\small
\begin{tabular}{@{}ll@{}}
\toprule
\textbf{Setting} & \textbf{Value} \\
\midrule
Starting checkpoints      & \pth{Helsinki-NLP/opus-mt-en-es}, \texttt{opus-mt-es-en} \\
Training pairs            & 1{,}998 per direction (training split only) \\
Development pairs         & 429 \\
Test pairs                & 432 (scored once) \\
Optimiser                 & AdamW \cite{loshchilov2019adamw}, learning rate $2\times10^{-5}$ \\
Batch size                & 8, gradient clipping at 1.0 \\
Maximum sequence length   & 128 tokens \\
Epochs run                & 4; checkpoint selected on development loss \\
Random seed               & 498 \\
Hardware                  & Apple M1 GPU via Metal Performance Shaders \\
Decoding (both systems)   & beam size 4, maximum 128 new tokens \\
\bottomrule
\end{tabular}\par}
\end{table}

The learning rate is deliberately small. A large rate would overwrite the general translation
knowledge acquired from OPUS, and $2\times10^{-5}$ is roughly an order of magnitude below typical
from-scratch rates, which is the standard range for adapting a pre-trained transformer. Only the
translation component was adapted; the reasoning for not fine-tuning the other components is given
in Section~\ref{sec:why_only_mt}.

\section{Dubbing Application}
\label{sec:webapp}

Objective~3 requires the synthesised dialogue to be integrated with the original video. The
research pipeline evaluates pre-segmented fragments, whereas real video contains continuous speech,
so a separate application was built around the same model instances --- it imports the pipeline's
model wrapper directly, so the application and the evaluation cannot drift apart. Its stages are
listed in Table~\ref{tab:webapp}.

\begin{table}[htbp]
\caption{Processing stages of the dubbing application. Stage~5 exists because of the measured
duration inflation reported in Section~\ref{sec:duration_result}.}
\label{tab:webapp}
\centering
{\singlespacing\small
\begin{tabular}{@{}clp{7.3cm}@{}}
\toprule
\textbf{Stage} & \textbf{Tool} & \textbf{Operation} \\
\midrule
1 & \texttt{ffmpeg}  & Extract the audio track to 16\,kHz mono PCM. \\
2 & Whisper          & Transcribe with segment timestamps, giving each utterance a slot on the timeline. \\
3 & MarianMT         & Translate each segment independently. \\
4 & XTTS-v2          & Synthesise each segment, conditioned on the full source audio for voice. \\
5 & WORLD            & Time-scale each synthesised segment to the duration of the slot it replaces. \\
6 & \texttt{ffmpeg}  & Remux: copy the video stream unchanged, replace the audio. \\
\bottomrule
\end{tabular}\par}
\end{table}

Stage~5 is the substantive one. Each synthesised segment is stretched or compressed to exactly the
duration of the segment it replaces using WORLD analysis and resynthesis \cite{morise2016world},
which changes duration without changing pitch, so formants --- and therefore the speaker's
identity --- survive the operation. Segments are then placed on a silent timeline at their original
start times, so the dubbed track has the same length as the source. Scaling is clamped to the range
$0.5\times$ to $2.0\times$, beyond which the audio degrades audibly; clamped segments are counted
and displayed rather than hidden, and Section~\ref{sec:dubbing_result} reports what happens when
the clamp binds. Stage~6 copies the video stream bit-for-bit, so the picture is never re-encoded.

\section{Computational Environment and Reproducibility}
\label{sec:reproducibility}

All experiments were run locally on an Apple silicon machine: the recogniser, synthesiser and
speaker encoder on the CPU, which is the stable configuration for this stack on that platform, and
translation fine-tuning on the integrated GPU through Metal Performance Shaders. No cloud compute
was used and no data left the machine.

Five practices support reproduction. Every experiment is driven by a command-line entry point
rather than by a notebook. Every result row is checkpointed to disk as it is produced, so a long
run can resume without repeating completed work. Every reported figure and table is regenerated
from the committed result files by a script, so no number in Chapter~\ref{ch:results} is typed by
hand. A provenance record is written beside the analysis outputs, capturing the git commit, the
interpreter version, the platform and the versions of every numerical library. And the exact
commands that reproduce each result are listed in Appendix~\ref{app:repro}.

\section{Ethical Considerations}
\label{sec:ethics}

No human participants were recruited and no personal data were collected for this study. The
corpus is a public-domain release whose own collection protocol obtained participant consent
\cite{ward2023dral}; it is used within its licence and is not redistributed. The pre-trained
models are used under their respective licences, and the synthesiser's non-commercial licence was
reviewed and accepted explicitly before use, through a deliberate opt-in rather than a default.

Voice cloning carries a dual-use risk that a dissertation on the subject should state plainly. A
system that reproduces a person's voice in a language they never spoke can be used to fabricate
speech attributed to them. Three properties of this work limit that risk without eliminating it:
the study operates only on a public-domain corpus whose participants consented to research use; no
model was trained to reproduce any specific individual, all cloning being zero-shot from a
reference recording; and the application produces no watermark-free artefact beyond what the
underlying public model already produces. Any deployment beyond research use should require
verified consent from the speaker whose voice is cloned, and should mark generated audio as
synthetic.
